\appendix
\appendixpage
\addappheadtotoc

\section{Additional material}
\label{sec:additional_fral}

\begin{proof}[Theorem~\ref{thm:fral_solving}]
3. $\imp$ 2. and 2. $\imp$ 1. hold by definition. By cyclic implications, it remains to prove 1. $\imp$ 3. Let $\la A, \var_A \ra$ be a $\mathcal T$-model over $X$. Suppose that $(FX, \var_{FX}) \models (X \vdash l = r)$. We want to prove that $A \brak{l} \var_A = A \brak{r} \var_A$.

The universal property of $FX$ gives a homomorphism $h : \underline{FX} \to \underline{A}$ such that\\ $h \circ \var_{FX} = \var_A$.
\begin{lemma}[Homomorphisms preserve term semantics]
  Let $A, B$ be two $\mathcal{T}$-algebras, $X \vdash t$ be a term and $h : A \to B$ be a $\mathcal{T}$-homomorphism. Then for every function $\var : X \to \underline{A}$, $h \bigl( A \brak{l} \var \bigr) = B\brak{r} (h \circ \var)$.
\end{lemma}
\begin{proof}
  This is a simple induction, since $h$ preserves the semantics of the operators in $\mathcal{T}$.
\end{proof}

\begin{lemma}[Homomorphisms preserve equations]
  Let $A, B$ be two $\mathcal{T}$-algebras, $\var : X \to \underline{A}$ be a function, $X \vdash l = r$ be an equation and $h : A \to B$ be a $\mathcal{T}$-homomorphism. Suppose that $A\brak{l} \var = A \brak{r} \var$. 
  
  \noindent Then $B \brak{l} (h \circ \var) = B\brak{r} (h \circ \var)$.
\end{lemma}
\begin{proof}
  \begin{align*}
    B \brak{l} (h \circ \var) &= h \bigl( A \brak{l} \var \bigr) &\text{since $h$ preserves term semantics}\\
    &= h \bigl( A \brak{r} \var \bigr) &\text{since $A \brak{l} \var = A \brak{r} \var$ by hypothesis}\\
    &= B \brak{r} (h \circ \var)  &\text{since $h$ preserves term semantics}\\
  \end{align*}
\end{proof}

Now we can finally prove our equality in $A$. By hypothesis, $FX\brak{l} \var_{FX} = FX\brak{r} \var_{FX}$. Since $h$ is a homomorphism and homomorphisms preserve equations, $A \brak{l} (h \circ \var_{FX}) = A \brak{r} (h \circ \var_{FX})$. But $h \circ \var_{FX} = \var_A$, therefore $A \brak{l} \var_A = A \brak{r} \var_A$.
\end{proof}

\begin{definition}[Abelian group monad]
	\label{def:abelian_group_monad}
	The Abelian group monad $\la A, \eta^A, \mu^A \ra$ is given by:
	\begin{itemize}
		\item $A : \Set \to \Set$ maps a set $X$ to the set containing all finite multisets with elements taken from $X$ and with multiplicities in the integers (i.e, elements of the multiset are allowed to have a negative multiplicity).
		\item $\eta_X^A : X \to AX$ maps an element $x \in X$ to the singleton $\{x : 1\}$.
		\item $\mu_X^A : AAX \to AX$ takes a union of multisets, accounting for all multiplicities. $$\mu_X^A\ \{\{a : 1, b : -2\} : 1, \{b : 1\} : 3\} = \{a : (1 \times 1), b : (-2 \times 1 + 1 \times 3)\} = \{a : 1, b : 1\}$$
	\end{itemize}
\end{definition}

\begin{definition}[Barycentric algebras]
	\label{def:barycentric_algebras}
	The theory of barycentric algebras $\BA$ consists of a binary operator $\oplus_p$ for every $p \in [0, 1]$ and the axioms \begin{align*}
    x &\vdash x \oplus_p x = x &\text{idempotence}\\
    x, y &\vdash x \oplus_p y = y \oplus_{1-p} x &\text{skew commutativity}\\
    x, y, z &\vdash x \oplus_p (y \oplus_r z) = \lp x \oplus_{\frac{p}{p+(1-p)r}} y \rp \oplus_{p + (1-p) r} z &\text{skew associativity}
  \end{align*}
\end{definition}

\begin{definition}[Join semi-lattices]
	\label{def:join_semilattices}
  The theory of join semi-lattices $\SL$ consists of a binary operator $\vee$ and a constant $\bot$, along with the axioms \begin{align*}
    x &\vdash x \vee \bot = x &\text{unit}\\
    x, y, z &\vdash (x \vee y) \vee z = x \vee (y \vee z) &\text{associativity}\\
    x, y &\vdash x \vee y = y \vee x &\text{commutativity}\\
    x &\vdash x \vee x = x &\text{idempotence}
  \end{align*}
\end{definition}

\begin{figure}
	\begin{mathpar}\mprset{flushleft}
(3 • x) • [2]\explain={⟨ Right neutrality ]}
(3 • x) • 2 • [ε]\explain={⟨ Left neutrality ]}
(3 • x) • 2 • ε • [ε]\explain={⟨ Left neutrality ]}
(3 • x) • 2 • ε • ε • [ε]\explain={⟨ Left neutrality ]}
(3 • [x]) • 2 • ε • ε • ε • ε\explain={⟨ Right neutrality ]}
(3 • [x • ε]) • 2 • ε • ε • ε • ε\explain={⟨ Right neutrality ]}
(3 • (x • ε) • [ε]) • 2 • ε • ε • ε • ε\explain={⟨ Left neutrality ]}
(3 • (x • ε) • ε • [ε]) • 2 • ε • ε • ε • ε\explain={⟨ Left neutrality ]}
(3 • (x • ε) • ε • ε • ε) • 2 • ε • ε • ε • ε\explain={⟨ Associativity ]}
3 • [((x • ε) • ε • ε • ε) • 2 • ε • ε • ε • ε]\explain={[ Associativity ⟩}
3 • [((x • ε) • ε • ε • ε) • 2] • ε • ε • ε • ε\explain={[ Commutativity ⟩}
3 • [(2 • (x • ε) • ε • ε • ε) • ε • ε • ε • ε]\explain={⟨ Associativity ]}
3 • 2 • ((x • ε) • ε • ε • ε) • ε • ε • ε • ε\explain={[ Associativity ⟩}
(3 • 2) • [((x • ε) • ε • ε • ε) • ε • ε • ε • ε]\explain={⟨ Associativity ]}
(3 • 2) • (x • ε) • [(ε • ε • ε) • ε • ε • ε • ε]\explain={[ Associativity ⟩}
(3 • 2) • (x • ε) • [(ε • ε • ε) • ε] • ε • ε • ε\explain={[ Commutativity ⟩}
(3 • 2) • (x • ε) • [(ε • ε • ε • ε) • ε • ε • ε]\explain={⟨ Associativity ]}
(3 • 2) • [(x • ε) • ε • (ε • ε • ε) • ε • ε • ε]\explain={[ Associativity ⟩}
(3 • 2) • ((x • ε) • ε) • [(ε • ε • ε) • ε • ε • ε]\explain={⟨ Associativity ]}
(3 • 2) • ((x • ε) • ε) • ε • [(ε • ε) • ε • ε • ε]\explain={[ Associativity ⟩}
(3 • 2) • ((x • ε) • ε) • ε • [(ε • ε) • ε] • ε • ε\explain={[ Commutativity ⟩}
(3 • 2) • ((x • ε) • ε) • ε • [(ε • ε • ε) • ε • ε]\explain={⟨ Associativity ]}
(3 • 2) • ((x • ε) • ε) • [ε • ε • (ε • ε) • ε • ε]\explain={[ Associativity ⟩}
(3 • 2) • ((x • ε) • ε) • (ε • ε) • (ε • ε) • [ε • ε]\explain={[ Left neutrality ⟩}
(3 • 2) • ((x • ε) • ε) • (ε • ε) • [ε • ε] • ε\explain={[ Left neutrality ⟩}
(3 • 2) • ((x • ε) • ε) • [ε • ε] • ε • ε\explain={[ Left neutrality ⟩}
(3 • 2) • [(x • ε) • ε] • ε • ε • ε\explain={⟨ Associativity ]}
(3 • 2) • (x • [ε • ε]) • ε • ε • ε\explain={[ Left neutrality ⟩}
[3 • 2] • (x • ε) • ε • ε • ε\explain={[ Evaluate ⟩}
5 • (x • ε) • ε • [ε • ε]\explain={[ Left neutrality ⟩}
5 • (x • ε) • [ε • ε]\explain={[ Left neutrality ⟩}
5 • [(x • ε) • ε]\explain={[ Right neutrality ⟩}
5 • [x • ε]\explain={[ Right neutrality ⟩}
5 • x
\end{mathpar}
\caption{\textsc{Frex}-extracted proof of $(3 • x) • 2 = 5 • x$ in the additive monoid over natural numbers.}
\label{fig:frex_proof_extracted}
\end{figure}
\newpage

% \section{Internship context}
% \label{sec:context}

% TODO

\section{Source code}
\label{sec:source_code}

The source code can be found on a fork of the \href{https://github.com/frex-project/idris-frex}{Idris Frex} project on GitHub at \url{https://github.com/S-c-r-a-t-c-h-y/idris-frex/tree/amaury}.

\section{Proof of Theorem~\ref{thm:correct}}
\label{sec:full_proof}

The full proof of Theorem~\ref{thm:correct} exceeds 50 pages and is compiled in a single file on GitHub at \url{https://github.com/S-c-r-a-t-c-h-y/semiring-solvers/blob/main/proof/semirings.pdf}.

\section{Research notes}
\label{sec:notes}

Research notes are available online for completeness on GitHub at \url{https://github.com/S-c-r-a-t-c-h-y/semiring-solvers/blob/main/notes/notes.pdf}.

\section*{Acknowledgments}

I would like to thank Ohad Kammar, my internship supervisor, for guiding me throughout my internship and helping me address any challenges I might have had.

I would like to thank Jesse Sigal for helping me with learning category theory, and Justus Matthiesen for helping me deal with some Idris 2 issues.

I would also like to thank Sergey Goncharov, Paul B. Levy, Cristina Matache, Sean K. Moss, Dominic P. Mulligan, Gordon D. Plotkin, John Power, Alex K. Simpson and Sam Staton on behalf of Ohad Kammar for fruitful discussions and suggestions prior to my internship.